\section{Related Work}
\label{sec:related}

\subsection{Underwater simulation and its fidelity}

Underwater robotics has converged on a small number of photorealistic,
physics-backed simulators, of which HoloOcean is the one used
here~\cite{potokar2022holoocea,potokar2024holoocea}; OceanSim and
comparable environments pursue the same goal of rendering and dynamics
close enough to support perception work~\cite{song2025oceansim}, and
recent digital-twin efforts push towards site-specific
replicas~\cite{mari2025digitalt,torroba2026marinari}. The value of such
tools rests on an assumption that is stated more often than it is
tested: that behaviour developed against simulated sensing survives
contact with real sensing. Work on simulation fidelity for autonomy has
begun to interrogate the assumption directly, both by quantifying how
sensing error propagates through a closed
loop~\cite{elmquist2022aperform,mahajan2024quantify} and by arguing that
sim-to-real evaluation protocols themselves need
rethinking~\cite{truong2022rethinki,kadian2020simrealp}. Our contribution
sits in that line, but differs in a way we consider important: rather
than measuring the gap after the fact, we measured the mismatches first,
built them into the simulator, committed the resulting predictions, and
only then went to the water.

\subsection{Monocular perception for underwater obstacles}

A camera is the sensor of choice on small inspection-class vehicles
because it is the sensor they already carry, and a great deal of work
addresses the resulting degradation: turbidity and colour
attenuation~\cite{yu2023udepthfa,drewsjr2016realtime}, learned monocular
depth in water~\cite{yang2022monocula,yang2025duvindif}, and
detection under the conditions a real basin
imposes~\cite{aldhaheri2025underwat,zhang2024localpat}. Benchmarks for
underwater perception have appeared~\cite{yang2025upbencha}, as have
uncertainty-aware treatments that acknowledge how unreliable the
resulting estimates are~\cite{lee2024uncertai,mustafa2023probabil}. What
this literature seldom reports is the quantity that mattered most in our
deployment: the range interval over which a given detector produces any
usable observation at all, given its own acceptance gates. We report it,
because in our case it -- rather than detection accuracy -- determined
the shape of the physical experiment.

\subsection{Reactive avoidance and its evaluation}

The dynamic-window approach remains the standard local
baseline~\cite{fox1997thedynam} and continues to be refined for marine
and dynamic settings~\cite{dobrevski2024dynamica,ma2025amodifie,%
eriksen2016amodifie}, alongside risk- and chance-constrained
formulations~\cite{han2025riskaware,han2026riskawar,mohamed2025chanceco}
and learning-based avoidance for underwater
vehicles~\cite{mari2026deeprein,li2024anobstac,wu2026collisio}. Vision-based
avoidance specifically for small ROVs has been demonstrated in pools and
tanks~\cite{alineipoiana2024abluerov,bergantin2024realtime,%
yao2023visionba}. We use a committed-manoeuvre planner and a holonomic
dynamic-window baseline, but we make no claim about their relative
merit in the water: only one of them was ever flown, for a reason we
report in Section~\ref{sec:deployment}, and the comparison between them
in this paper is entirely simulated.

\subsection{Sim-to-real transfer reported honestly}

Sim-to-real transfer for marine vehicles is an active
topic~\cite{zhang2025simtoreal,groot2025scenario,meyers2025testinga},
and the field is increasingly explicit that a transfer result is only as
strong as the protocol that produced it. Pre-registration is rare in
robotics; where it appears, it is usually as a fixed evaluation protocol
rather than as predictions committed before the physical trial. Our
design takes the stronger form, and then encounters the case that makes
the distinction worth having: the physical system that was eventually
flown was not the system that had been frozen. We report that outcome
rather than repairing it retrospectively, and we separate every
post-deployment simulation from the frozen set at the level of notation.

We claim no new avoidance algorithm, simulator or detector: the
contribution is the study design of Section~\ref{sec:introduction} and
what it exposed. We make no priority claim for the combination, only the
observation that auditing one's own frozen campaign, and stating which
physical quantities a camera-only campaign cannot measure, are not
common practice in this literature.
